\section{Evaluation}
\label{sec:evaluation}

\subsection{Experimental Setup}
\label{subsec:setup}

Using the pipeline from Section~\ref{sec:methodology}, we analysed the Rust standard library for the host target \texttt{aarch64-apple-darwin}. The corpus consists of the \texttt{core}, \texttt{alloc}, and \texttt{std} crates compiled with monomorphisation enabled, so each concrete instantiation of a generic function contributes its own MIR body and CFG. In total, the corpus contains 27{,}997 functions: 22{,}273 safe functions and 5{,}724 functions declared as \texttt{unsafe fn}.

For each function, we computed the treewidth of its MIR CFG. We report the mean, median, and maximum treewidth, together with the share of functions with treewidth at most~3 and at most~6. The threshold of~6 is included to facilitate comparison with the bound established for goto-free C and for label-free Java~\cite{thorup_1998,gustedt_2002}.

\subsection{Results}
\label{subsec:results}

Treewidth is very small across the corpus. Over all 27{,}997 functions, the mean treewidth is 1.04, the median is 1, and the maximum observed treewidth is 6. Moreover, 99.65\% of all functions have treewidth at most~3, and all functions have treewidth at most~6. Thus, MIR CFGs in the Rust standard library are overwhelmingly close to trees in practice.

\begin{table}[h]
	\centering
	\begin{tabular}{lrrrrr}
		\hline
		Group            & Count    & Mean & Median & Max & $\mathrm{tw} \leq 3$ \\
		\hline
		All functions    & 27{,}997 & 1.04 & 1      & 6   & 99.65\%              \\
		Safe functions   & 22{,}273 & 1.06 & 1      & 6   & 99.59\%              \\
		Unsafe functions & 5{,}724  & 0.97 & 1      & 4   & 99.90\%              \\
		\hline
	\end{tabular}
	\caption{Treewidth summary for the Rust standard library MIR corpus.}
	\label{tab:evaluation-summary}
\end{table}

\noindent Table~\ref{tab:evaluation-distribution} shows that the distribution is strongly concentrated at very small values: 97.90\% of all functions have treewidth at most~2, and only 97 functions (0.35\%) exceed treewidth~3. The largest values occur in a small number of formatting and systems-oriented routines with unusually branch-heavy control flow.
\begin{table}[h]
	\centering
	\begin{tabular}{lrr}
		\hline
		Treewidth & Count  & Share \\
		\hline
		0         & 3{,}853 & 13.76\% \\
		1         & 19{,}773 & 70.63\% \\
		2         & 3{,}783 & 13.51\% \\
		3         & 491    & 1.75\% \\
		4         & 90     & 0.32\% \\
		5         & 5      & 0.02\% \\
		6         & 2      & 0.01\% \\
		\hline
	\end{tabular}
    \caption{Distribution of treewidth values across all functions in the Rust standard library
   MIR corpus.}
    \label{tab:evaluation-distribution}
\end{table}

The safe/unsafe split does not indicate higher structural complexity for unsafe Rust in this corpus. Unsafe functions have a slightly lower mean treewidth than safe functions (0.97 versus 1.06), and their maximum observed treewidth is only~4.

\subsection{Discussion}
\label{subsec:discussion}

Overall the results support the hypothesis that practical Rust CFGs have small treewidth. Even without a Rust-specific structural proof analogous to Thorup's result~\cite{thorup_1998}, the empirical picture is clear: almost all MIR CFGs lie well below the bound of~6 associated with structured C and Java fragments, and the maximum observed value is exactly~6.

For further context, the results invite comparison with the Java study of Gustedt et al., which should be interpreted with care: their analysis is source-level and structurally constrained, while ours is based on MIR basic-block CFGs~\cite{gustedt_2002}.
In particular, treewidth~0 and~1 are common in our setting but cannot occur in their construction due to Thorup's decomposition algorithm.
With this in mind, they report an average treewidth of about~2.7 and a maximum of~5 for the Java standard library, whereas we obtain a mean of~1.04 and a maximum of~6 for Rust MIR.

While these results are promising, several limitations should be considered.
First, the study is limited to the Rust standard library and may not represent application code or third-party crates.
Second, the analysis is performed on optimised MIR for a single compiler version and host target, so the exact CFGs may vary across releases and platforms.
Third, we analyse undirected CFGs after removing unwind-only cleanup edges, which matches the graph-theoretic notion of treewidth but abstracts away some aspects of Rust's exceptional control flow.

Despite these limitations, the results provide a consistent empirical baseline for the structural complexity of Rust MIR CFGs.
The consistently small treewidth values suggest that compiler optimisations designed for graphs of bounded treewidth are well-suited to real Rust MIR in practice.
